本研究使用了生成式人工智能工具，并在下表中按环节如实记录其产出与相应的人工核验方式。所有结论性判断、建模决策与最终取舍均由作者作出并逐项核对；AI 未被用于伪造数据或结果。全部数据处理、数值计算、测试、绘图与排版在 Modal 云平台执行，本地只进行文本编辑、版本管理与产物散列核对。

\begin{longtable}{@{}p{0.14\linewidth}p{0.40\linewidth}p{0.40\linewidth}@{}}
\toprule 环节 & AI 产出 & 人工核验 \\* \midrule \endhead \bottomrule \endfoot
问题理解与\newline 方案比较 &
阅读题目与附件说明，整理题目未规定之处（边界条件形式、附件 1 之外的烘房状态、收缩域上方程的写法、结果文件时间范围），列出各自的备选方案及其后果 &
逐条核对题目原文与附录公式（含 $D$ 中温度的单位）；对每一处解释撰写建模决策记录（\nolinkurl{docs/mdr/}），并要求程序同时给出备选结果以供对照 \\
建模 &
由守恒律推导一维径向热质耦合方程与 Robin 边界条件；给出物质坐标（Landau）下收缩域方程的推导与守恒性证明；提出离散守恒与极值原理两个命题 &
人工复核推导的每一步与量纲一致性；独立验算物质坐标下 $R$ 因子的相消；命题的证明由作者复核后写入论文 \\
程序编写 &
编写物性与边界数据模块、径向有限体积求解器（分段自适应 BDF、事件求根、稀疏雅可比）、与求解器分离的校验器（贝塞尔级数解、守恒、极值原理、二维轴对称交叉检验）、收敛与灵敏度实验、结果工作簿与图表生成程序及其单元/性质测试 &
审阅全部源代码（论文附录“完整程序”收录）；核对贝塞尔级数解的系数与特征值方程；核对结果工作簿与附件 3 模板的逐格一致性；所有数值结论以校验器报告与收敛表为准 \\
实验组织 &
在 Modal 上组织全部阶段的运行，整理结果说明（\nolinkurl{docs/RESULTS.md}）与数据说明（\nolinkurl{docs/DATA_NOTES.md}） &
核对每个阶段的清单（输入输出散列、参数、环境）；抽查结果文件与论文表格的一致性；确认论文中的每个数字都来自阶段登记的数值宏或生成表格 \\
论文写作 &
按竞赛论文格式组织章节、撰写与改写文字、绘制技术路线图与图表说明 &
逐节审阅并修改表述；核对全部图表被正文引用、公式被引用、参考文献真实可查；核对摘要与正文数字与结果文件一致 \\
\end{longtable}

\section*{附：工具与运行环境}
\begin{itemize}
  \item 生成式 AI：Anthropic Claude Code（模型 Claude Fable 5.1），用于上表所列环节。
  \item 云计算平台：Modal，用于全部数据处理、数值计算、单元测试、图表生成、\LaTeX{} 排版与打包；镜像内的依赖版本固定并记录在每个阶段的清单中。
  \item 数值库：SciPy（\texttt{solve\_ivp} 的 BDF/Radau、\texttt{PchipInterpolator}、Bessel 函数与求根）、NumPy、pandas、Matplotlib、openpyxl。
\end{itemize}
